% chapters/ch11.tex —— 第 11 章
% 体例：具体失败 → 设计问题 → 机制解释 → 对照实验 → 条件变化 → 选择原则
% 数字身份纪律：
%   - 2601.22628（TTCS）、2506.01716（自出题）、2505.03335（零数据自博弈）、
%     2508.03680（训练框架）的数字一律 \cited{}；
%   - 本书三次运行的读数一律 \measured{}；
%   - 由上述数字算出的倍数与差值一律 \calcnum{}。
% E07 / E09 未运行；需训练资源与冻结协议，不只依赖黑盒API余额。
\chapter[经验数据、自生成课程与 Agent 强化学习]{经验数据、自生成课程\\与 Agent 强化学习}
\label{ch11}

\chaptergoal{下一轮学什么由谁决定，多步交互里的反馈又怎样训到执行策略上？}
\chapterprops{本章为P2（怎样区分反馈与更新对象的贡献）
补上另一半观测位：当任务与奖励都由系统自己产生时，反馈质量的含义变了。
E07与E09未运行，本章据外部研究和教学推导讨论课程机制。}
\chapterfig{图~\ref{fig:ch11-loop}：自生成课程怎样用求解反馈调整下一轮出题，以及在哪里单列外部评价}
\chapterpitfall{训练分数一路在涨，回头一看题库越来越软——没有一个外部校验点在上面数着}

% ---------- 1. 具体失败 ----------
\begin{failurebox}
TTCS 让合成器出题、求解器解题，两者从同一个预训练模型初始化，交替训练。
以 Qwen2.5-Math-1.5B 为起点，六项数学基准的均值从 \cited{17.30} 提高到 \cited{41.49}。

作者分别去掉合成器训练和重复度惩罚，两种设置在 AMC23 上都从 \cited{62.50} 降到 \cited{55.00}。论文给出两种解释。冻结合成器后，求解器的能力仍在变化，生成的题可能太易或太难。保留合成器训练而去掉重复度惩罚，出题者则倾向于反复改写相似的问题。要辨别这两条机制，可以跟踪各阶段题目对当期求解器的难度分布，以及同批题目中有多少道只是换了说法\cite[§5.4及表3]{ttcs-2601}。

这两组对照引出本章的问题：当系统自己生成训练题时，
若只看不断变化的自生成训练题上的分数，就无法单独分出\textbf{是能力涨了，还是题库变软了}。上面引用的基准评价仍是相应适配条件下的效果证据，不能因训练信号来自模型内部就将它一并否定。
\end{failurebox}

\designquestion{上一章决定了经验写在哪个位置；本章决定\textbf{数据流从哪来}。
这里有两个前后相接的问题，必须分开回答：一是\textbf{学什么}——下一轮练哪一批任务；
二是\textbf{怎么学}——多步交互里的奖励与信用怎么分派。
若同时改变任务与学习方式，最终差异就难以归因；失败也可能表现为分数下降或训练不稳定，并非必然上涨。}

前一半（\S\ref{sec:ch11-data}、\S\ref{sec:ch11-filter}、\S\ref{sec:ch11-synth}）
处理数据：成功与失败轨迹、筛选与去重、任务生成与环境验证。
后一半（\S\ref{sec:ch11-mdp}、\S\ref{sec:ch11-reward}）处理学习：
马尔可夫决策过程的建模、信用分配与奖励投机。
后续处理工程对接与长期运行（\S\ref{sec:ch11-plumbing}、\S\ref{sec:ch11-nonstationary}）。

% ---------- 2. 成功与失败 ----------
\section{成功轨迹与失败轨迹的不同价值}
\label{sec:ch11-data}

\begin{sloppypar}
先给不含糊的那一半。一组自出题工作让同一个 8B 模型只用自己合成的数据训练，
四个环境的均值 pass@1 从零样本的 \cited{12.0} 提到 \cited{23.5}（自改进设置），
约 \calcnum{2.0} 倍；若把更强的模型（70B）在这批合成任务上的示范轨迹拿来当监督信号，
还能再走到 \cited{32.2}（蒸馏设置）（\cite{self-challenging-2506} §5.2，Table 1；
\cited{800} 条合成任务、\cited{12k} 条离线 rollout，pass@1 为四次独立试验均值）。
\end{sloppypar}

\begin{sloppypar}
失败轨迹在这两种设置里的处理不同，涉及生成策略、反馈与训练目标：
\end{sloppypar}

\begin{itemize}
\sloppy
  \item \textbf{师生之间}。蒸馏设置下，原作者的初步结果显示，
        弱模型也可能从强模型的失败轨迹中受益，作者举出的可学信息包括对错误的反思；
        他们因此在该设置下不分成败地用全部轨迹训练（\cite{self-challenging-2506} §4.3 式 (3)）。
  \item \textbf{自我迭代}。自改进设置下奖励是 $0/1$ 的结果信号，
        固定已生成任务与验证函数时，在同策略、无基线的一步奖励梯度中，零奖励项不贡献梯度；与成功轨迹微调的联系还受归一化等条件约束（同前引 §4.3 式 (2)，详见第10章）。
\end{itemize}

\begin{sloppypar}
失败是否值得保留取决于它提供什么信息。错误操作、失败原因与后续修正可以具有诊断价值，但不应把错误动作直接当成正确示范。原论文关于教师失败反思的观察有具体设置，不能推成“只有更强教师才能提供有用失败数据”；执行反馈也可以使自行探索的失败具有可检验含义。
\end{sloppypar}

% ---------- 3. 筛选 ----------
\section{任务生成、环境验证与自博弈}
\label{sec:ch11-synth}

% figures/ch11-loop.tex —— 第 11 章核心图（全书引用入口）
% 用法：在 book/chapters/ch11.tex 中 % figures/ch11-loop.tex —— 第 11 章核心图（全书引用入口）
% 用法：在 book/chapters/ch11.tex 中 % figures/ch11-loop.tex —— 第 11 章核心图（全书引用入口）
% 用法：在 book/chapters/ch11.tex 中 \input{figures/ch11-loop}
% 图主体在 ch11-loop-body.tex，便于单独编译与复用。
\begin{figure}[htbp]
\centering
\input{figures/ch11-loop-body}
\caption{自生成课程的反馈循环。沿上方路径，生成的任务经筛选后交给求解器，表现或一致性信号反馈给下一轮出题。向下的箭头单列评价与保持，读图时应区分参与筛选的任务与留作独立评价的任务。}
\label{fig:ch11-loop}
\end{figure}

% 图主体在 ch11-loop-body.tex，便于单独编译与复用。
\begin{figure}[htbp]
\centering
% 分析框架，非所有方法共同实现。
\begin{tikzpicture}[
 b/.style={draw,rounded corners=2pt,text width=3.2cm,minimum height=1.2cm,align=center,font=\small},
 a/.style={-{Stealth[length=2mm]},thick}]
 \node[b,fill=orange!10] (gen) at (0,0) {任务生成\\来源、变体与版本};
 \node[b,fill=yellow!10] (filter) at (4.8,0) {课程筛选\\执行检查、去重\\或内部难度代理};
 \node[b,fill=blue!8] (solve) at (9.6,0) {求解与学习\\轨迹、反馈与策略};
 \draw[a] (gen.east)--node[above,font=\scriptsize]{候选}(filter.west);
 \draw[a] (filter.east)--node[above,font=\scriptsize]{保留}(solve.west);
 \draw[a] (solve.north) -- ++(0,0.65) -| node[pos=.25,above,font=\scriptsize]{表现或一致性代理反馈}(gen.north);
 \node[b,fill=green!8] (eval) at (9.6,-2.3) {评价与保持\\明确范围与独立性};
 \draw[a] (solve.south)--(eval.north);
 \node[align=left,text width=7cm,font=\small] at (2.4,-2.3)
 {筛选通过不保证正确；多数一致不代表真值。\\评价题参与选择后，应登记其开发用途。};
\end{tikzpicture}

\caption{自生成课程的反馈循环。沿上方路径，生成的任务经筛选后交给求解器，表现或一致性信号反馈给下一轮出题。向下的箭头单列评价与保持，读图时应区分参与筛选的任务与留作独立评价的任务。}
\label{fig:ch11-loop}
\end{figure}

% 图主体在 ch11-loop-body.tex，便于单独编译与复用。
\begin{figure}[htbp]
\centering
% 分析框架，非所有方法共同实现。
\begin{tikzpicture}[
 b/.style={draw,rounded corners=2pt,text width=3.2cm,minimum height=1.2cm,align=center,font=\small},
 a/.style={-{Stealth[length=2mm]},thick}]
 \node[b,fill=orange!10] (gen) at (0,0) {任务生成\\来源、变体与版本};
 \node[b,fill=yellow!10] (filter) at (4.8,0) {课程筛选\\执行检查、去重\\或内部难度代理};
 \node[b,fill=blue!8] (solve) at (9.6,0) {求解与学习\\轨迹、反馈与策略};
 \draw[a] (gen.east)--node[above,font=\scriptsize]{候选}(filter.west);
 \draw[a] (filter.east)--node[above,font=\scriptsize]{保留}(solve.west);
 \draw[a] (solve.north) -- ++(0,0.65) -| node[pos=.25,above,font=\scriptsize]{表现或一致性代理反馈}(gen.north);
 \node[b,fill=green!8] (eval) at (9.6,-2.3) {评价与保持\\明确范围与独立性};
 \draw[a] (solve.south)--(eval.north);
 \node[align=left,text width=7cm,font=\small] at (2.4,-2.3)
 {筛选通过不保证正确；多数一致不代表真值。\\评价题参与选择后，应登记其开发用途。};
\end{tikzpicture}

\caption{自生成课程的反馈循环。沿上方路径，生成的任务经筛选后交给求解器，表现或一致性信号反馈给下一轮出题。向下的箭头单列评价与保持，读图时应区分参与筛选的任务与留作独立评价的任务。}
\label{fig:ch11-loop}
\end{figure}


\begin{sloppypar}
自出题这条路要成立，先得解决一件事：这到底是不是一道题。
一种可检查的做法是把部分要求写成代码。一组自出题工作把每个生成的任务定义成四件套——
指令、验证函数、示例解、显式失败案例——除指令外全部用代码表达，
交给执行器检查示例能否通过、指定失败案例能否被拒绝；这些检查只能排除部分无效题
（\cite{self-challenging-2506} §1、§4.2）。
\end{sloppypar}

\begin{sloppypar}
零数据自博弈那条路把这一招推到极致：同一个模型既出题又解题，
代码执行器统一负责验证任务是否合法、判定答案是否正确，
执行结果成为其编程任务的可核对反馈（\cite{absolute-zero-2505} §1、§3）。执行器独立于语言生成不等于任务语义被完整证明；测试覆盖、资源限制与题目构造仍需核对。
\end{sloppypar}

\begin{table}[htbp]
\centering
\caption{自生成课程在三个基座上的成绩（六项数学基准的均值；原论文 Table 1）}
\label{tab:ch11-ttcs}
\begin{small}
\begin{tabular}{>{\raggedright\arraybackslash}p{3.6cm}ccc}
\toprule
方法 & Qwen2.5-Math-1.5B & Qwen2.5-Math-7B & Qwen3-4B-Base \\
\midrule
预训练基座 & \cited{17.30} & \cited{26.02} & \cited{34.58} \\
自一致性（多数投票） & \cited{27.62} & \cited{32.15} & \cited{42.07} \\
二分角色共进化 & \cited{31.70} & \cited{42.07} & \cited{39.98} \\
测试时强化学习 & \cited{36.56} & \cited{48.42} & \cited{43.59} \\
共进化加课程合成 & \cited{41.49} & \cited{52.54} & \cited{47.21} \\
\bottomrule
\end{tabular}
\end{small}
\end{table}

\begin{sloppypar}
读这张表要注意两件事。第一，收益的量级不小（1.5B 那一列从 \cited{17.30} 到 \cited{41.49}，
差 \calcnum{+24.19}），但它是\textbf{在同一批测试题上做测试时适配得到的成绩}，
与“模型在开放任务上变强了”不是一回事；本书只把它当参照。
第二，按原表均值计算，三种基座相对各自起点的增量分别为\calcnum{24.19}、\calcnum{26.52}与\calcnum{12.63}个平均分数点；7B 的增量比1.5B多\calcnum{2.33}，不是两者最终成绩差\calcnum{11.05}。这些不同模型的结果不构成只改变模型规模的受控实验。原文 Figure 3 另有跨基准与一般领域评价，不能将主表的测试时适配边界误写成全文完全没有迁移证据。
\end{sloppypar}

\subsection{生成一道题，需要同时检查哪几个对象}
\label{sec:ch11-taskcontract}

一个自然语言要求、一段验证代码和一个示例答案能够互相不一致。示例通过验证，只说明这两段程序在某次执行中相容，不保证它们满足文字要求。Self-Challenging 的 CaT 将指令、验证函数、示例解与失败案例一起生成并执行检查；其任务质量分析恰好展示了这种不一致的具体后果\cite{self-challenging-2506}。

原论文在 Retail 环境分析中报告，只有验证函数的条件经过可执行过滤后保留\cited{47.7\%}任务；再要求示例解通过，保留比例为\cited{9.5\%}；加入失败案例拒绝检查的完整 CaT 为\cited{5.2\%}。作者为每个条件取得\cited{50}条任务求解轨迹，由作者人工分类\cite{self-challenging-2506}。过滤保留比例与人工误判比例是两种不同分母，不能把\cited{5.2\%}读成错误率，也不能把人工复核冒称本书做过的检查。

作者发现一个值得保留的反例：仅要求示例解通过，可能放大过于宽松的验证函数。最宽松的验证器什么也不做就判成功，当然能让示例通过，却没有确认任务完成。加入明确失败案例，可以排除这类漏洞。在报告的人工分析中，完整条件没有出现相应假阳性，但仍有假阴性；其中一个问题是任务指令缺少必要信息\cite{self-challenging-2506}。有限样本中没有观察到一种错误，不能扩成验证器永远不会出现这种错误。

将这个机制写成自己的任务合同，可以分别核对四种关系。指令要给足完成任务所需的条件；示例解应在规定环境中完成指令；验证函数应接受满足要求的结果；失败案例应覆盖几种明确违反要求的结果。这四种关系不能由同一个“运行成功”标记概括。尤其要区分程序没有异常、示例返回成功，以及任务目标真的完成。

考虑一个教学任务：“从列表中返回满足条件的所有元素，并保持输入顺序。”若验证器只检查返回元素均满足条件，它会接受空列表，漏掉“所有”；若只检查集合相等，它可能接受乱序或重复结果。示例代码在一个正常输入上通过，并不能发现这两类问题。加入遗漏元素、乱序、重复和空输入等反例，才会覆盖不同要求。这里的反例来自明确任务语义，是本书教学设计，不是现有实验台的新任务集。

失败案例也可能过于贴合当前示例。若每次只把正确输出的最后一个字符改错，验证器只需识别这个表面变化就能通过检查。可以按违反的要求组织反例，再让具体输入变化：缺少必要动作、范围超出、顺序错误、使用过时状态或错误对象。这样做不会穷尽所有失败，但有助于说明目前覆盖了什么，而不是用一个反例数量充当质量证明。

另一个常被漏掉的对象是环境状态。示例解在带有某条记录的数据库中可行，不代表求解器运行时那条记录仍存在。如果验证器只检查最终状态，还可能把任务开始前已经满足的条件当成智能体完成了目标。应记录初始状态与允许的变更，区分“环境本来正确”与“系统作出了所要求的动作”。这对于有副作用的任务尤其重要。

执行器本身也要有范围。代码能运行不代表可以任意访问机器状态；只读检查、允许写入的对象、超时与输出上限应由执行环境约束。将生成代码交给解释器，只建立了执行通道，没有自动建立隔离。当前本书教学实验的具体保护边界见第14章和附录B，本章不把课程设计建议写成已具备的安全能力。

\section{多数答案一致，不等于多数答案正确}
\label{sec:ch11-consensus}

TTCS 在两个环节使用求解器对同一道题的多次回答。合成器依据回答的一致程度估计题目难度，求解器则以回答中出现最多的答案作为训练用的伪标签\cite{ttcs-2601}。先看这批回答能告诉我们什么。

设同一题采样得到 $K\geq1$ 个答案，频数最高的答案为 $\hat y$（并列时按预定规则处理），一致率为
\begin{equation}
 s(x)=\frac{1}{K}\sum_{k=1}^{K}\mathbf1[y_k=\hat y].
 \label{eq:ch11-consistency}
\end{equation}
这个式子没有调用标准答案。若所有回答都落在同一个错误答案上，$s(x)$ 仍等于一；若回答包含许多不同错误，一致率可能很低，但低一致也不能证明题目有足够信息或确实可解。它测量的是指定采样与答案归一化条件下的集中程度，正确性需要另一条证据。

用一个教学例说明：十个回答中六个给出甲、四个给出乙。多数一致率为\calcnum{0.6}。若甲正确，正确率同为\calcnum{0.6}；若乙正确，正确率是\calcnum{0.4}；若二者都错误，正确率为\calcnum{0}。观察同一份答案频数，存在三种不同正确性情形。因而不能把一致率直接命名为通过率，也不能仅凭一致率接近中间值就断言题目处于真实能力边界。

论文用 $R_{\mathrm{cap}}=[4s(1-s)]^\gamma$ 表示题目对当前求解器的难度是否适中。题目通过格式检查后，从这个分数扣除两项惩罚，一项衡量它与原题的相似程度，另一项衡量同批生成题的重复程度。设两项加权之和为 $R_{\mathrm{sim}}$，最终奖励就是 $\max(0,R_{\mathrm{cap}}-R_{\mathrm{sim}})$。格式检查失败时奖励为零。论文取\cited{$\gamma=1.2$}，求解器在训练时保留满足 $|s-0.5|\leq\delta$ 的样本，取\cited{$\delta=0.25$}\cite{ttcs-2601}。

作者公布的代码给出了另一种扣分次序。若一道变体被判为与原题过于相似，先将 $R_{\mathrm{cap}}$ 乘以 $1-r$（$r$ 为触发判定的相似度），再扣除组内重复惩罚，结果最低为零\cite{ttcs-code-e29e95c}。同样的 $r$ 在这里使高分题损失更多；在论文公式里，相同的参考题惩罚对应相同的扣分量，直至触及零下限。论文两个版本的式~(8)--(9) 都采用后一种写法。主表数字仍按论文给出的条件解读，复算单题奖励则需注明采用哪个版本。

还要区分三组采样数：合成器为一道输入题生成\cited{4}个变体，能力评估对每个变体采样\cited{10}次回答；求解器训练通常为每题采样\cited{8}条轨迹，AIME24、AIME25 与 AMC23 则使用\cited{16}条\cite{ttcs-2601}。出题数、评估回答数与训练轨迹数承担不同用途，不能都用前面的十次回答例子代替。

仍在十次回答的教学条件下，$s=\calcnum{0.6}$ 的未取幂项为\calcnum{0.96}，而 $s=\calcnum{0.9}$ 的对应项为\calcnum{0.36}。这个算式准确反映代理偏好，却没有使用哪一个答案真正正确的信息。若较难题的错误答案也高度集中，它可能被当作过于一致的样本排除；若不完整题目导致回答分散，它可能落入保留区间。二者都属于需要检查的代理失配，不是原论文已测出的发生率。

一致率还依赖答案提取。数学表达可以存在不同但等价的形式，代码可以有多种正确实现；若只按字符串相等分组，语义相同的回答可能被拆散。反过来，过度归一化可能把不同答案合并。应保存原输出、提取结果和归一化规则，才能判断课程变化来自模型行为还是解析器变化。这个要求对外部正确性检查同样成立。

即使额外假定独立采样，也要区分“与预先固定的正确答案一致”和“与这批回答选出的多数答案一致”。考虑只有甲、乙两种答案的教学模型：甲正确，每次独立答甲的概率为 $p$，十次中答甲的数量 $Z\sim\operatorname{Binomial}(10,p)$。样本正确率为 $Z/10$，多数一致率却是 $\max(Z,10-Z)/10$。当 $p=\calcnum{0.6}$ 时，前者的期望为\calcnum{0.6}，后者按二项分布求和得到约\calcnum{0.6470}。多数标签由同一批样本选出，因此一致率并非正确率的无偏估计。这个算例只检验估计量的含义，不是 TTCS 的实测误差。

若标签预先固定且确为正确答案，独立正确性指示变量的均值才有相应的无偏性与集中界。原文附录C.1以正确性伯努利变量解释课程代理\cite{ttcs-2601}，将这种解释用于无标签多数奖励时，需要另查多数标签的可靠性和样本依赖。共享模型、提示与训练来源的回答还可能共享错误，增加次数可以使错误多数更稳定；次数本身不能提供独立真值。

训练时，一致性提供了一种无标签奖励。评价时，研究者用基准答案检查求解器能否解对题。TTCS 的主表与迁移实验回答的是后一个问题；要检验多数答案作为伪标签有多可靠，还需比较每批回答选出的答案与金标。

\subsection{课程的难度应相对于谁、在哪个时点定义}
\label{sec:ch11-difficulty}

一道任务对旧模型很难，对更新后模型可能很容易；对会使用某个工具的执行器容易，对没有该工具的执行器又可能不可解。因此课程表里的“难度”至少要绑定求解器版本、提示和工具条件。脱离这些条件的单一难度标签，只能是材料属性或人的判断，不等于当前系统的成功概率。

在线估计难度会消耗求解机会。用当前模型试做候选，再根据结果决定是否进入训练，这批试做就是课程选择成本。若把这些试做也用于更新，应明确它同时承担两个用途；若丢弃，也不能从资源账中删除。题目被挑选时使用的模型与真正训练时的模型相隔很久，还可能产生难度估计过期的问题。

一种朴素的课程策略是偏好既非全对也非全错的任务。在真值奖励下，这可能提供有差异的组内学习信号；在一致性奖励下，它偏好的是答案分布不完全集中。两种情况不能共用“有一定通过率”的解释。即便使用真值，所有回答失败也不意味着任务永远不可学：示范、工具支持或更细的子任务可能改变条件。课程只是当前训练预算下的选择，不是对任务价值的永久裁决。

还有一种风险是把偶然成功当成跨过能力门槛。少量采样中出现一次正确解，可能源于猜测或特殊路径；少量全败也可能来自波动。可保留每次试做的输出和反馈，而不只保存一个难度均值。选择规则应在运行前说明怎样处理样本不足、缺失反馈和并列候选，不要在看见结果后用有利规则补齐。

课程应当追求什么也需要先定义。最短时间内提高一组目标题，广泛掌握新任务家族，以及维持旧能力同时学习新任务，是不同目标。第一种可以合理地围绕目标题生成变体；第二种需要检查跨来源迁移；第三种还需安排保持评价和回放。把三者都称为自生成课程有效，会让看似冲突的研究结果失去比较基础。

本书建议用一张任务家族乘时间的记录表表达课程：每轮生成了哪些家族、保留多少、训练多少、旧锚点表现如何。表格不是越细越好，而要能回答“下一轮为何选这批题”。如果选择理由只是当轮代理分数最高，应如实记录；不要在事后给每次选择补写一个模型当时并未使用的学习理论解释。


% ---------- 5. MDP 与信用分配 ----------
\section{数据筛选、去重与难度控制}
\label{sec:ch11-filter}

\begin{sloppypar}
不同研究使用不同的筛选信号。Self-Challenging 通过执行生成的验证函数、示例解和失败案例检查部分一致性；TTCS 则用多数答案一致率构造难度代理，并在奖励中加入相似性惩罚。后者不是外部标准答案正确率。
同一组实验把后两个部件逐个拆掉，读数如下：
\end{sloppypar}

\begin{table}[htbp]
\centering
\caption{拆掉三个部件分别掉多少（原论文消融表，本书只转述，不做差值之外的推断）}
\label{tab:ch11-filter}
\begin{small}
\setlength{\tabcolsep}{4pt}
\begin{tabular}{>{\raggedright\arraybackslash}p{4.0cm}ccc}
\toprule
设置 & AMC23 & Olympiad & Minerva \\
\midrule
完整设置 & \cited{62.50} & \cited{36.05} & \cited{40.44} \\
去掉合成器训练 & \cited{55.00}（\cited{$-7.50$}） & \cited{32.79}（\cited{$-3.26$}） & \cited{37.50}（\cited{$-2.94$}） \\
去掉在线数据过滤 & \cited{60.00}（\cited{$-2.50$}） & \cited{33.68}（\cited{$-2.37$}） & \cited{38.60}（\cited{$-1.84$}） \\
去掉多样性惩罚 & \cited{55.00}（\cited{$-7.50$}） & \cited{35.31}（\cited{$-0.74$}） & \cited{39.34}（\cited{$-1.10$}） \\
\bottomrule
\end{tabular}
\end{small}
\end{table}

\begin{sloppypar}
三行是分别移除一个部件后的点估计，支持比较这些具体消融条件，不足以单独识别中间机制。移除合成器训练、移除在线过滤、移除多样性惩罚都可能同时改变保留数据量和训练分布；若要解释损失为何发生，还需记录重复、可行性、曝光与预算。原文把部分下降归因于重复或难度不匹配，本书保留其解释身份，不把表格改写成已经验证的因果图。
\end{sloppypar}

\begin{pitfallbox}
\textbf{反例：静态的强出题者打不过会变的弱出题者。}
同一组实验把一个更强的模型（14B 指令模型）当成固定的合成器，原文称其余框架保持不变，
在Table~2的AIME24、AIME25与Minerva三项汇总中，相对同一参照的提升为\cited{+2.66}；共进化条件的提升为\cited{+5.34}
（\cite{ttcs-2601} §5.3 Q3，Table 2）。这里同时改变了合成器模型身份与是否训练，支持该组合对照下共进化条件的点值更高，不能独立证明可适应性普遍比模型强度重要。Table 3 同一小模型的冻结消融补充了训练因素证据，但仍不等于完整的模型强度乘训练方式析因实验。
\end{pitfallbox}

% ---------- 4. 任务生成 ----------
\subsection{过滤会改变课程，不只是清除坏题}
\label{sec:ch11-selection}

先生成再过滤，看起来只是删去不合格数据，但保留后的分布可能与生成分布很不一样。容易执行的任务、短示例和模板清楚的题目，可能比复杂任务更容易通过同一组检查。因此，过滤质量与保留覆盖需要分别观察；保留率越低并不自动表示课程越好，也不自动表示课程崩溃。

Self-Challenging 还按示例解长度比较过滤前后的难度代理分布，观察到较弱生成器的分布被收窄，并与较强生成器作比较\cite{self-challenging-2506}。这是一种分布诊断，但示例长度不是能力难度的唯一量尺。短程序可以要求关键洞察，长程序也可能重复简单操作。本书不将其结果概括成过滤后所有难度保持不变。

设计自己的课程账时，可以给每个候选保留固定标识、生成来源、预计类别和最终去向。去向至少区分格式不合、执行失败、示例未通过、反例未被拒绝、重复、难度代理不合以及真正进入训练。若多个条件同时失败，要说明只记首个原因还是记全部原因；否则不同筛选顺序会改变每个原因的计数。

例如，一道题先被判重复后就不执行，系统不知道它是否也会执行失败。把未执行题全部当成执行通过会高估质量，把它们全部当成执行失败也不合理。串行过滤适合节省成本，但各项通过率要注明分母。为了估计被早期过滤遮住的问题，可以对少量候选安排完整检查；这是额外审计设计，不能从现有首失败日志里事后恢复不存在的结果。

重复也需要多种观察。完全相同的输入是最容易识别的一类；替换实体名称而保留同一求解结构，是另一类；不同表面结构却要求同一关键技能，又可能是有价值的多样练习。把所有相似性都惩罚掉，会排除巩固同一技能所需的变体；只去掉字面重复，则可能留下大量模板复述。应先说明希望覆盖的技能差异，再选择近似检测方法。

课程质量还受任务家族权重影响。假如某个家族容易合成并通过检查，按自然产量取样会逐步增加它的训练权重。总体训练成功率随之上升，可能只是容易家族占比提高。可以同时报告按实际训练权重的成绩和固定家族权重的成绩，前者描述系统真实经历，后者提供跨轮比较的参照。两者无需相同，但各自的分母要稳定可查。

成本方面，被拒绝的候选仍然付出生成和检查费用。若每个保留任务来自许多被丢弃的尝试，只按训练样本数报告数据效率会漏掉主要支出。应同时记录候选数、保留数、实际训练曝光、模型调用与执行检查量。原文的保留比例可以说明为什么这项成本值得关注，却不能在没有运行条件的情况下直接推算本书生成一个任务要多少钱。

\subsection{一个可手算的课程漂移案例}
\label{sec:ch11-driftexample}

以下数据完全是教学设定，用于检验一种解释，不来自模型运行。设课程只有甲、乙两个家族。旧版本在甲上的成功率为\calcnum{0.90}、在乙上为\calcnum{0.40}；新版本分别变成\calcnum{0.80}和\calcnum{0.30}。两个家族都下降了\calcnum{0.10}，但如果训练题的构成同时改变，总分可以朝相反方向移动。

假定旧课程的甲占\calcnum{10\%}、乙占\calcnum{90\%}，旧版本的加权成功率为\calcnum{0.45}。新课程将甲提高到\calcnum{90\%}、乙降低到\calcnum{10\%}，新版本在新课程上的加权成功率为\calcnum{0.75}。若每轮只报告当轮课程均值，曲线就上升\calcnum{0.30}；但固定任一家族观察，能力都下降了。

这个反例不声称真实自生成系统总会这样变化，它证明“变化题集上的均值上涨”不足以单独推出各家族能力提高。若用固定各占一半的权重评价，旧版本为\calcnum{0.65}、新版本为\calcnum{0.55}，下降就重新可见。实际报告可以同时保留自然权重与固定权重：前者反映正在服务或训练什么，后者帮助追踪可比范围内的变化。

还可以做版本与题集的交叉评价。旧模型在新课程上的成功率为\calcnum{0.85}，新模型在旧课程上为\calcnum{0.35}。同一个新题集上，新模型仍比旧模型低\calcnum{0.10}；同一个旧题集上也是如此。四格比较把“课程变容易”与“模型在相同范围变差”区分开，但在真实实验中，需要保存两份题集和两个可运行版本，不能只靠两次总体均值重建。

若某个家族后期完全消失，仅有自然课程日志就没有它的新版本表现。缺失格不能补成零，也不能沿用旧成绩；应该安排固定评价或承认这部分未观测。生成器可能出于合理业务变化减少该家族，也可能在奖励驱动下避开难题。原因需要进一步记录与对照，单凭分布变化还不能判定动机。

这个例子也说明为什么“课程覆盖率”需要外部分母。若把当轮出现的家族当作全部家族，每轮覆盖都能是百分之百。应预先定义目标范围，或者明确采用开放探索的覆盖描述而不声称穷尽。新发现的家族可以加入索引，但旧版与新版覆盖分数的分母变化必须注明。

\subsection{一道筛选习题：分母怎样从日志中消失}

再看一份教学流水账。生成器产生\calcnum{100}个候选，先删除\calcnum{20}个字面重复；剩下\calcnum{80}个中，\calcnum{16}个格式不合；对余下\calcnum{64}个执行示例检查，其中\calcnum{32}个通过；最后有\calcnum{8}个通过指定失败案例的拒绝检查并进入训练。数字是人为设置的，目的只是训练读账能力。

最终保留率相对全部候选为\calcnum{8\%}，相对进入最后一道检查的候选为\calcnum{25\%}。两者都可以报告，但若只写“通过率百分之二十五”，读者会不知道前面的候选去了哪里。示例检查通过率是\calcnum{50\%}，其分母是实际执行的六十四个，不能把它写成所有候选中有一半可解。

更不能据此说另外九十二个都是无效题。二十个重复候选可能完全可解，只是没有新信息；十六个格式不合候选可能经过修正后可用，但本次没有这样运行；后面的失败可能来自题目、验证器或环境。日志记录的是规则之下的去向，任务语义质量还需要额外证据。把所有拒绝原因合并成“错误”，会损失改进生成器所需的信息。

如果下一轮最终保留率从\calcnum{8\%}升到\calcnum{16\%}，至少存在几种解释：生成质量改善、题目变简单、去重规则放宽、验证器变宽松，或者候选类别变了。只有固定检查版本并核对各阶段分母，才有条件逐步排除解释。即使最终确认保留质量提高，也还需要检验训练这些题目是否带来目标能力收益。

读者可以尝试重新排序过滤步骤。把便宜的格式检查提前，可能减少后续执行；先做去重，则可能避免重复计算。这是资源优化问题，但重排后每项“首失败”计数会变，不能把变化本身当成生成器质量改善。若所有检查确定且彼此无副作用，最终接受集合可以相同；若检查会修改状态、调用随机模型或依赖前一结果，这个条件就需要实际验证。

同一份账还应记录重试。一个候选执行三次才取得有效反馈，仍是一个候选，却发生三次检查尝试。按候选计课程规模，按尝试计资源，二者各有用途。将重试展开成三个新任务会夸大数据多样性；将三次只记成一次费用则低估开销。这个小例子对应附录F的唯一事件账原则，在课程场景中同样适用。

\section{状态、动作、观察与结束条件}
\label{sec:ch11-mdp}

\begin{sloppypar}
前面各章可以把一次执行当成一次函数调用；一旦要做强化学习，就得把话说清楚。
一份训练框架工作给出的表述可以当作本书的模板：
区分完整执行状态与模型可见观察，原文 §3.2 明确用部分可观测马尔可夫决策过程表述；一次模型输出作为动作，环境执行后转移并提供相应奖励（\cite{agent-lightning-2508} §3.1—3.2）。
在这个表述下，一条轨迹就是一串转移，每次模型调用都能被单独拿去训练，
于是多轮交互里最麻烦的两件事有了落脚点：累积上下文让序列越来越长，
以及多智能体场景下“哪一步该为结果负责”。
\end{sloppypar}

\begin{itemize}
\sloppy
  \item \textbf{信用分配}。做法是在层级化的强化学习里加一个信用分配模块，
        把回报分配给模型调用，再接入单轮算法；原文当前实现给同一回合的每个动作相同终局回报，并未学出逐动作因果贡献（同前引 §3.3，Figure 3）。
  \item \textbf{中间奖励}。结果奖励太稀疏时，可以按系统监控信号给中间奖励，
        例如工具调用的返回状态成功与否（该文称之为自动中间奖励，同前引 §3.4）。
        这增加了反馈密度，也要求检查工具返回成功是否对应任务进展，下一节具体讨论。
  \item \textbf{结束条件}。任务最多允许跑多少步，是参与信用分配的超参数，
        不是无关紧要的常数：改变上限可能改变被截断轨迹和反馈机会；原本已提前结束的轨迹则未必变化。
\end{itemize}

% ---------- 6. 奖励投机 ----------
\subsection{一条轨迹进入训练前，还缺什么}
\label{sec:ch11-transition}

执行日志通常为排查故障而写，训练数据却必须说明模型在何种条件下作出何种动作。日志里有最终答案和总奖励，仍可能缺少每次调用的原始输入、实际输出、策略版本及结束原因。把这类日志直接称为完整轨迹，会在后续更新时引入无法追溯的空白。

Agent Lightning 将执行状态组织为语义变量，模型调用只能看到其中一部分；其 §3.2 因而明确采用部分可观测的表述\cite{agent-lightning-2508}。状态是环境和执行器当前的完整描述，观察是实际交给策略的输入。若工具返回被摘要后才进入下一次调用，训练时应保留实际摘要输入，而不能事后用完整返回替换它，否则训练条件与行为产生时不同。

以检索问答为例，一次调用产生查询，工具返回候选文档，下一次调用产生答案。查询动作和答案动作属于同一回合，却面对不同观察。最终答错可能来自查询偏离、检索器未召回、文档过时或回答误读。只保存最后一轮，会把前面造成的约束消失掉；保存了全部调用，也不意味着已经知道每一步的因果贡献。

训练接口还应区分模型生成内容与工具提供内容。工具返回是下一步观察的一部分，不是当前策略选择输出的词序列。若把整段拼接文本都作为策略动作训练，可能错误地奖励或惩罚模型并未生成的内容。具体算法可以使用掩码、按调用拆分或其他表示，但需要验证实际参与损失的范围，而不只检查文本能否重新加载。

结束原因也是学习信号的一部分。任务完成、超出步数、环境异常、服务拒绝和反馈缺失，不能一律折成“策略失败”。某些协议会把超限按失败计入主要任务指标，这是可以预先规定的评价规则；但训练与故障诊断仍应保留原始原因。否则基础设施不稳定可能被学成对某类任务的错误偏好，或者被无意过滤出结果表。

策略版本需要绑定到动作。并行执行中，旧版本开始的回合可能在新版本发布后才结束；若中途路由到新版本，整条轨迹还可能混合多个行为策略。训练算法对旧数据的适用条件应由算法说明，不能把所有刚到达的数据都称为同策略数据。记录时间戳有帮助，但真正需要的是可以核对的模型或适配器版本与采样配置。

将这些字段收齐，是为了让学习问题可定义，不是为了给日志增加装饰。一个最小的可审计样本应能回答：模型看到了什么，输出了什么，哪个环境执行了它，得到了什么反馈，之后为何继续或结束。额外需要的概率、优势或工具状态取决于算法。缺少某个必需字段时，应标记不适用，而不是用默认值伪装成真实观测。

字段齐全之后，成功轨迹还要适合目标学习器。Agentic-Teacher 使用外部编码教师，从视觉—语言—动作模型的失败状态构造修复代码策略；给成功策略的机械臂动作加入时间相关噪声后重新执行，保留成功轨迹，再删除闲置帧、保留夹爪切换，并按冻结基座的流匹配损失剔除损失异常高的轨迹，用于后训练\cite{agentic-teacher-2611}。这里的损失衡量与该基座学习目标的兼容程度，不是真值判据。论文在仿真中研究这条流程；其累计处理对照不能分离每个部件的独立贡献。

\subsection{把终局奖励分给动作，不等于识别动作贡献}
\label{sec:ch11-credit}

Agent Lightning 的当前实现将同一回合中的各次模型动作赋予相同终局回报，再交给单轮学习算法处理词元层更新\cite{agent-lightning-2508}。这提供了一个简洁、可用的训练接口，但其含义不是“已证明每一步同样有用”。相同赋值是算法选择，因果贡献是另一个需要证据的问题。

考虑一条教学轨迹：先作出一个偏离目标的查询，发现无关结果后重新检索，最后正确回答。终局奖励为一时，统一回报会把两个查询和最终回答都与成功相联。这个估计仍可能在大量样本与适当基线下提供学习信号，但单条轨迹不能告诉我们第一个查询是必要探索、可以省掉的错误，还是偶然无害的绕路。

若为了“更精确”而给每个动作人工分数，也需说明分数依据。把工具调用成功作为正反馈，可能奖励语法正确却没有帮助的查询；把错误返回作为负反馈，可能惩罚有价值的探测。所谓过程奖励必须对应任务目标或有明确用途的代理，不能因为发生在中途就默认比终局奖励更真实。

可以用反事实复评研究动作贡献，但它也有条件。将某一步替换为另一动作后重跑下游，改变的不仅是该步输出，还包括后续观察、模型随机性和环境状态。需要从可复现快照开始，控制允许变化的部分，并说明下游是否重新采样。若环境无法重置，就不能把两个不同时间的回合作为严格反事实。反事实方法增加的成本也应计入研究预算。

归一化会进一步改变动作权重。假设两条教学轨迹都成功，一条包含两个模型调用，另一条包含六个。若拆成调用后对全部调用等权，后一条贡献的样本数量为前者的\calcnum{3}倍；若先在回合内平均再对回合平均，两条轨迹的总权重可以相同。两种目标并非同一个实现细节，应按希望优化的单位选择。词元长度平均又会引入第三层权重。

长轨迹拥有更多训练项，不代表它更有价值；短轨迹也不一定更优，因为任务难度可能不同。可以在日志中同时保留回合、调用和词元三个层级，明确损失在哪一级求和、在哪一级平均。第10章成功样本数与全样本数的归一化问题，在这里变成多层版本；任何“算法等价”的主张都需对齐这些层级。

组内相对奖励也需要有差异。若一个组的所有样本得到完全相同奖励，减去组均值后，这一奖励项的优势为零；加入数值稳定项不会凭空产生目标方向，另加的 KL 正则项仍可能产生梯度。课程选择可能因此偏好奖励不全相同的题目。若差异来自解析错误或随机验证噪声，训练仍可能朝错误方向走。

但原始奖励方差较大，并不表示标准化后的优势方差也较大。设一组二元奖励的均值为 $s$，按组内总体标准差归一化，暂略数值稳定项；当 $0<s<1$ 时，
\begin{equation}
 A_i=\frac{r_i-s}{\sqrt{s(1-s)}},\qquad
 \frac1K\sum_i A_i=0,\qquad \frac1K\sum_i A_i^2=1.
 \label{eq:ch11-normalized-advantage}
\end{equation}
十个样本中六个奖励为一、四个为零，原始方差是\calcnum{0.24}；九个为一、一个为零，原始方差是\calcnum{0.09}。两组标准化后的组内方差却都为\calcnum{1}，后一组的唯一零奖励对应优势\calcnum{$-3$}。因此，不能只从 $s(1-s)$ 推出实际策略梯度更大；还要看优势与得分函数梯度的配对、裁剪、稳定项和正则项。原文附录C.1的方差论证\cite{ttcs-2601}在应用时，也须明确讨论的是哪一种量。

可以把信用分配理解为一种有假设的学习规则，而不是审计系统自动给出的责任认定。训练规则要以最终效果评价；调试时的责任分析要结合执行证据与必要的干预。这两个用途可以相互帮助，却不能互相替代。将它们分开，也能避免把一句“采用层级强化学习”写成已经解决多步归因。

\section{可验证奖励与奖励投机}
\label{sec:ch11-reward}

\begin{sloppypar}
把反馈写成可执行规则有助于复查，但规则仍可能不完整或有实现错误。可重复计算与真实目标一致是不同要求。
该框架在检索问答任务上将词级答案F1与格式项组合，奖励写成
$R=\cited{0.9}\times R_{\text{答案F1}}+\cited{0.1}\times R_{\text{格式}}$
（\cite{agent-lightning-2508} §4.2）。格式项可以单独提高：
模型按要求输出标签即可提高格式项，即使正确性未变，总奖励也可能上升。格式本身可以是合法需求；这里展示的是复合奖励的解释风险，不是原论文已经观察到奖励投机的证据。要判定投机，应另查真实目标是否未改善或受损，以及行为如何利用了指标缺口。
\end{sloppypar}

\begin{sloppypar}
还需要结合任务上下文解释训练中出现的输出。Absolute Zero 报告了一段令作者担忧的推理文本：Llama-3.1-8B 在第\cited{132}步构造任务时，讨论怎样设计难解、复杂的程序。其上游任务提议提示本就要求算法难度与复杂度，不能将这段文本转述为模型在没有指令背景时自行设立了目标。
作者将其作为安全关切的例子，并说明安全管理自我改进组件不在本文处理范围内
（\cite{absolute-zero-2505} §6、Figures 34 与 36）。
该记录支持分析生成行为与提示的关系；它没有展示实际执行危害，也不能直接代表所有自博弈系统。
\end{sloppypar}

\begin{sloppypar}
本章应区分执行检查、模型生成的评审和内部一致性信号。执行检查不等于形式证明，在线过滤也未必使用真值。E02 的标签扰动设置及其保留执行线索的限制见第4章；它不能替代本章尚未运行的自生成课程对照。
\end{sloppypar}

% ---------- 7. 执行与训练 ----------
\subsection{复合奖励上升时，先看哪一项在变}
\label{sec:ch11-rewardledger}

原文检索问答实验的奖励包含词级答案F1与格式两项，权重分别为\cited{0.9}和\cited{0.1}\cite{agent-lightning-2508}。奖励拆分不是旁注，而是读训练曲线的必要信息。一个总分只能显示加权和，不能指出模型改善了答案、格式，还是两者之一掩盖了另一项下降。

设教学例中平均答案F1保持\calcnum{0.4}，格式分数从\calcnum{0}升到\calcnum{1}。总奖励从\calcnum{0.36}升到\calcnum{0.46}，增加\calcnum{0.10}。这完全符合给定公式，却没有发生正确性改善。它并不自动构成违规行为：若结构化输出是服务需求，格式提高本身有价值。问题在于不能将这条总奖励曲线命名为回答能力曲线。

再设格式分增加一，而答案F1从\calcnum{0.4}下降到\calcnum{0.3}。总奖励由\calcnum{0.36}变为\calcnum{0.37}，仍增加\calcnum{0.01}。这展示了加权和可能允许的交换。是否接受这种交换，是目标设计问题；如果业务要求正确性不得下降，就需要单独约束，而不能事后指责模型遵从了预先给定的标量目标。

奖励投机的判断需要更具体。可以事先定义任务完成指标及不可接受行为，再检查高奖励是否来自利用验证缺口，而不是完成要求。例如任务要求检索证据后回答，系统却输出形式合法但不受证据支持的内容；此时要有独立的证据支持检查，才能确认它偏离了目标。只观察总分上升和某个案例失败，不能估出整个训练过程的投机率。

监控表可以同时保留奖励各分量、最终任务指标、资源消耗和违规事件。这样，当总分上升时，可以追问究竟哪一项推动了变化。若某些分量只在训练任务上可得，而最终评价使用另一检查器，也应说明差异。这能帮助区分代理过拟合与真实迁移失败，避免把所有不一致都归为模型“作弊”。

修改奖励后，旧曲线与新曲线通常不能直接连起来。权重变了、解析器变了或某项从缺失变成可测，都会改变尺度。可以用新规则离线重算保存完整的旧轨迹，但重算结果应标为重新评分，不是旧训练当时使用的奖励。若旧记录不够重算，应保留断点，而不是拼出一条看似持续提高的曲线。

\section{执行框架与训练框架的连接}
\label{sec:ch11-plumbing}

\begin{sloppypar}
从跑一个智能体到更新一次参数，中间隔着一整套工程：
轨迹要落成训练样本，更新后的模型要回到执行侧继续被用。
一份框架工作给出的答案是\textbf{把两者彻底拆开}：
训练侧保留一个服务端，管理训练过程，并向外暴露一个与 OpenAI 形式兼容的接口；
执行侧跑一个客户端，一边运行智能体，一边把轨迹送回去
（训练与执行解耦，\cite{agent-lightning-2508} §3.4，Figure 4）。
这种接口减少执行框架与训练实现的耦合；具体系统仍需处理轨迹、奖励、版本与工具适配，不能保证接入任何现有系统都无需修改。
\end{sloppypar}

\begin{table}[htbp]
\centering
\caption{三个验证任务：不同框架、不同工具、被优化的执行者也不同}
\label{tab:ch11-plumb}
\begin{small}
\begin{tabular}{>{\raggedright\arraybackslash}p{2.8cm}>{\raggedright\arraybackslash}p{2.8cm}>{\raggedright\arraybackslash}p{2.4cm}>{\raggedright\arraybackslash}p{2.2cm}>{\raggedright\arraybackslash}p{3.4cm}}
\toprule
任务 & 框架 & 数据集 & 工具 & 被优化的部分 \\
\midrule
文本转 SQL & LangChain & Spider & SQL 执行器 & 三个角色用同一模型；选其中两类转移训练 \\
\midrule
开放域问答 & OpenAI Agents SDK & MuSiQue & 维基检索器 & 单模型、单条指令 \\
\midrule
数学问答 & AutoGen & Calc-X & 计算器 & 单模型，含工具调用 \\
\bottomrule
\end{tabular}
\end{small}
\end{table}

\begin{sloppypar}
有一件设计值得单独写下来：\textbf{可以选择性地优化多智能体系统里的某一部分}。
第一个任务选择“写 SQL”与“改写 SQL”角色的转移用于训练，不选择检查角色的转移。这不自动等于检查角色所调用的参数或行为被冻结。
这对第 9 章的角色结构演化是一个直接可用的接口：
结构搜索可以决定有哪些角色，训练侧可以选择哪些转移参与更新，两者不必同步。若角色共用被更新的权重，未入选角色的输出仍可能变化；要声称它被冻结，还需记录独立权重副本、适配器或路由边界。
\end{sloppypar}

% ---------- 8. 非平稳 ----------
\section{非平稳任务、回放与能力遗忘}
\label{sec:ch11-nonstationary}

\begin{sloppypar}
自生成课程让训练任务分布与当前模型相互影响。旧任务遗忘指标并不因此失效：只要保留固定旧任务集和明确的评价版本，仍可计算保持与后向迁移。真正失去参照的是每轮只评价新生成题目、却把不同题集的平均分当作同一能力曲线的做法。
\end{sloppypar}

\begin{itemize}
\sloppy
  \item \textbf{回放与复评}。把早期合成题重新用于训练叫回放，固定旧题定期重考则是复评。前者改变学习分布，后者提供保持参照，不能混用名称。
        本书没有跑过带回放与不带回放的对照，所以给不出它的代价。
  \item \textbf{参数侧同时在忘}。第 10 章那个连续更新后部分早期任务表现受损的结果
        （\cite{seal-2506} §5，Figure 6）说明参数这一侧确实会忘；
        本章的结果说明课程这一侧会变形。二者可能相互作用，不能未经模型与数据支持就写成简单相加；固定锚点与版本交叉评价有助于拆分。
  \item \textbf{三类需分别检查的问题}：固定出题者与求解器的匹配、合成题的复述与重复，
        以及自博弈中令作者担忧的生成行为。E07 为前两类课程问题准备对照：
        比较固定课程、随机任务、自生成课程三者在\textbf{独立迁移集}上的差别，
        需事先定义足够有用的迁移收益；训练上涨而迁移差异不显著并不足以反对该主张，还要检查区间能否排除预定幅度。
\end{itemize}

\begin{sloppypar}
E09 研究交互式强化学习相对离线轨迹蒸馏的增量，但重新采样会改变数据，因此需分别定义固定数据的目标比较与允许重新交互的系统比较。区间包含零不能单独证伪有用性。按当前蓝图，两项均为教学级设计，尚缺冻结协议、可训练模型和计算资源；黑盒API可用不能解除这些条件。
\end{sloppypar}

\subsection{固定锚点、旧题回放与新题探索}
\label{sec:ch11-replaydesign}

自生成训练分布可以不断变化，评价仍可保留固定锚点。锚点是跨版本重复测量的参照，回放则是重新进入训练的旧材料。两者的功能不同：锚点帮助发现变化，回放尝试影响变化。把用于训练的旧题重新考一遍，可以观察拟合保持，却不能再称为完全未参与训练的独立评价。

可以把材料按用途维护为三个集合：允许反复训练的旧经验、用于开发选择的回归题，以及保留给最终评价的旧能力题。它们可以共享能力类别，但不应在文件管理中失去身份。如果根据某道回归题失败而专门生成相似练习，该题所代表的信息已经参与后续设计，应在来源关系中记录，而不是只看原题文本是否被复制。

新题探索与旧题回放会竞争预算。增加回放可以降低新材料曝光；只保留近期任务则可能忽视早期能力。没有一条固定比例适合所有系统。本书建议把比例视为开发条件，记录新旧家族的实际曝光及对应保持指标，再在限定资源内比较。这里没有本书的回放对照实测，因此不推荐一个看似精确的通用比例。

还应区分遗忘与目标变化。如果业务要求已经废弃某条旧规则，继续遵守它未必是需要保留的能力；如果仍需回答历史版本问题，则应把时间或版本条件写入题目。把所有旧回答变化都计为遗忘，会惩罚必要更新；把所有变化都称为适应，又可能掩盖真正损失。锚点集需要记录为什么该能力仍然有效。

评价矩阵可以将模型版本放在行、固定任务家族放在列，并另列各版本实际训练分布。前者显示行为变化，后者帮助提出解释。仅凭两张表的相关变化仍不能确定因果，但至少可以发现“只在新题上测量”造成的参照漂移。第~\ref{sec:ch10-retentionmatrix}~节给出另一种汇总方式。它按材料在序列中的位置对齐多次运行，列内会包含不同任务。要追查某一道题从哪个版本开始退化，仍须保留任务标识。

停止条件应覆盖学习与数据两个层面。若固定锚点持续受损，可以停止发布新候选；若生成题的有效覆盖明显收窄，可以暂停扩大课程；若预算耗尽，应结束当前试验并保存未完成状态。停止不是承认自生成方法无效，而是对当前设计作可复查的资源和质量决策。重新启动时，需说明改变了什么条件，不能把中断前后的轨迹无标记地拼接。


\section{怎样拆开课程收益与学习算法收益}
\label{sec:ch11-factorial}

E07 关注任务如何产生，E09 关注数据怎样用于学习。若一次对照同时更换课程生成器、训练算法、工具和推理预算，最终可以评价整套系统，却难以知道下一次该保留哪个部件。教学设计应先决定主要问题，再选择能够回答它的比较，而不是靠实验编号给所有差异一个统一解释。

固定数据比较可以从同一份带来源的轨迹包开始，让两种更新目标使用相同可用信息，并限定开发与训练预算。这有利于研究目标和优化的差异，但无法完整代表需要在线交互的算法。允许重新采样时，算法会遇到新状态、获得新反馈，比较就包含探索与数据获取效果；应将这部分费用和信息增量一起报告。

课程比较同样有两种方式。冻结候选题池，再比较固定顺序、随机抽取与能力驱动选择，可以较清楚地研究选择规则；让生成器各自在线创造题目，则更接近完整自生成系统，但题目内容与选择规则会一起变化。前者不能证明开放生成有效，后者不能单独归因于难度排序。两个问题都值得研究，名称应与实际操作一致。

若资源允许，可以做一个小型析因设计：一个因素是课程来源，另一个因素是学习方式，同时固定基座和评价范围。它有助于观察某种课程是否只对某种算法有用。但增加条件并不自动带来足够判决力，尤其当每个条件只有极少独立任务家族或训练种子时。当前本书未冻结这种设计，不能把它当作已经安排好的确认性实验。

关于“同预算”，至少要列出数据生成、候选试做、过滤、训练、开发选择与最终评价的边界。固定最终训练步数只能控制其中一项。某种方法训练较快却需要昂贵题目筛选，另一种方法训练较慢却直接使用固定题库，总费用可能与训练时间排名不同。预算不足导致少完成训练时，应按预定规则报告中止和已用资源，不悄悄删去该条件。

最终评价还需匹配主张。若声称适配目标题，应报告适配时允许看到哪些题目和标签；若声称迁移，应在未用于生成与选择的范围评价；若声称长期能力积累，应同时检查固定旧任务与新任务。不能在结果出来后挑一种最有利的评价身份，再将其他未改善的范围移到脚注。

E09 的结果解释尤其要避免“区间跨零即两种学习一样”。需要先规定什么增量有实际价值，再根据差值与不确定范围判断是支持有用优势、排除足够大的优势，还是尚无法区分。小样本的宽区间可能同时容纳很大收益与损失，最诚实的结论就是需要更多证据。改变算法名称或加长说明不能补足判决力。

\subsection{从本章方法到一份最小实验记录}
\label{sec:ch11-minrecord}

准备一次课程试验时，最小记录不应从“用了哪种算法”开始，而应从一次完整数据流开始：候选任务怎样产生，为什么被留下，谁尝试了解答，得到何种反馈，哪些片段进入更新，哪个版本接受了评价。每个环节都应能回到实际工件。只保留最终模型和平均分，无法检查前面讨论的大部分问题。

可以先做不调用模型的流程检查。例如，用确定的教学候选验证重复检测、缺失反馈处理、训练与评价标识隔离、预算扣减和中止保存。这样的离线检查能确认程序按约定处理数据，却不能证明生成器会产生好任务或求解器会进步。将流程检查记成实验成功，是把软件行为与学习效果混为一谈。

真实运行前还需固定哪些修改允许发生。生成器可以改题目，不代表可以改最终评分标准；策略可以学习工具选择，不代表可以提高自己的调用额度；课程可以调整新旧数据比例，不代表可以悄悄删除旧能力评价。这些边界属于研究设计，不因系统具有自修改能力而自动放开。

试验结束时，至少对账候选数量、入选数量、实际完成的回合、成功取得反馈的回合、进入训练的样本和参与最终评价的对象。各项不一定相等，但差额都应有去向。中断运行可以提供故障诊断和部分描述，不应补写未执行对象的结果，也不能把不同设计的半截运行合并成一个完整对照。

当前本书 E07 与 E09 尚未完成这些前置条件。本章的设计与教学计算使后续运行可以更明确地规划，同时保留实测缺口。这比给未运行的路线写一个预期提升值更有用：将来无论结果正、负或不确定，都有一套事先说明的语言记录它。


% ---------- 9. 条件变化 ----------
\subsection{条件变化：自生成课程在什么时候会骗人}
\label{sec:ch11-conditions}

\begin{itemize}
\sloppy
  \item \textbf{验证器与生成器共享假设时}，应检查共同错误。模型审查可以作初筛，但其结论不能自动充当独立正确性证明；外部执行器也只覆盖编码进去的要求。
  \item \textbf{没有独立迁移评价时}，结论应限制在适配或训练范围。训练指标是过程证据，不能单独支持未见任务上的迁移主张；测试时适配可有自己的明确研究目标。
  \item \textbf{出题者与解题者同源时}，要额外监视两类退化：
        一是题库退化（多样性惩罚挡的就是这个），
        二是 Uh-oh Moment 那一类作者也没预料到的行为。
        这两类问题都不能仅由一条总体奖励曲线可靠诊断，需要任务和行为记录。
  \item \textbf{单题噪声大时}，任何“某道题变差了”的判断都要先报多次运行的结果。
        本书在自己的任务集上观察到 \measured{16} 道题里有 \measured{9} 到 \measured{12} 道
        出现多次运行结果不一致（第 3 章）。
\end{itemize}

% ---------- 10. 能说与不能说 ----------
\section{读三组论文结果时，不要混用评价身份}
\label{sec:ch11-resultidentity}

本章引用的三条路线可以相互启发，却不能直接按一个最终百分数排名。TTCS 在无标签测试题及其变体上适配，Self-Challenging 生成可执行任务并训练求解器，Agent Lightning 主要提供执行与训练连接方式。它们改变的数据、计算过程和评价范围不同；先确定这些身份，才能提出有意义的比较。

\begin{sloppypar}
TTCS 在每项基准上分别进行测试时训练。两项 AIME 的\cited{mean@32}是在温度\cited{0.6}下，对每题\cited{32}次回答的正确性取平均，再按题汇总；它不表示“至少答对一次”的 pass@32，也不是多数投票准确率。其他数学基准使用温度\cited{0}的贪心解码，数学答案的格式等价性判断还由 GPT-4o-mini 辅助处理（\cite{ttcs-2601} 附录A.2）。主表六项均值汇总了这些不同评价条件。测试时训练允许利用待评价题目本身及其变体，但不利用其真值来适配，这与参数完全冻结的零样本评价不同。
\end{sloppypar}

原文 Figure 3 另外展示在指定数学基准上适配后，对其他数学基准以及一般领域题集的评价\cite{ttcs-2601}。这为跨范围表现提供了额外证据，不能因主表采用测试时训练就把全文描述成没有迁移分析。反过来，这些已报告范围也不等于所有开放任务上的持续提升。正文应逐项说清适配集、选择过程和评价集，而不是在“泄漏”与“普遍智能”两个极端标签之间选择。

Self-Challenging 主表在四个工具环境报告 pass@1 和 pass@4，并说明生成\cited{800}个合成任务和\cited{12k}条离线轨迹，pass@1 为\cited{4}次试验平均\cite{self-challenging-2506}。多次作答的成功机会与单次成绩不是同一个指标；将 pass@4 与另一方法的单次成功率并排，必须同步交代尝试预算。正文引用的\cited{12.0}、\cited{23.5}和\cited{32.2}来自同表相应 pass@1 汇总，不能从另一列挑更好数字替换。

其中蒸馏与自改进的差别也不只在算法名称。蒸馏条件使用更强模型产生轨迹，自改进条件使用自己的轨迹；若直接相减，差异包含数据生成者。论文另有在线算法比较，仍需看其独立设置，而不能用主表师生差直接估计“强化学习额外收益”。本书 E09 必须重新定义对照，不能借用这三个数替代未运行的实验。

Agent Lightning 的任务包括文本转SQL、开放域问答与工具数学问答，所优化角色及奖励各有差别\cite{agent-lightning-2508}。执行接口能够连接多种工作流，是工程范围证据；它不意味着同一个奖励构造对所有工作流都正确。评价框架支持一个接口，与某个新任务已经获得可靠训练效果，是两种不同层次的完成状态。

读者可以在每张结果表旁写下六个字段：训练材料来源、行为模型、可见反馈、适配时点、评价指标与资源边界。若一个字段未知，就保留未知。这样做比跨论文直接计算一个综合排名更有帮助，因为它会指出下一项受控比较需要补什么，也能防止把来源不同的数字拼成一条未发生过的训练曲线。

\section{让课程、学习规则与评价各有参照}
\label{sec:ch11-claims}

自生成课程把系统的选择扩展到训练材料本身。下一轮题目是否成立、反馈是否对准任务、当前模型能否从中得到有效信号，是三个相接但不同的问题。可执行任务合同、内部一致性代理与独立范围评价分别提供部分证据，应按各自作用解释。

设计E07时，应区分固定候选池上的选择规则与在线创造新题的整体流程；设计E09时，应区分固定数据的更新目标与允许重新交互的系统。两项实验均未运行。已有论文为具体机制提供了观察，本书的教学算例帮助检查分母、权重与参照，并不填补训练结果。

\falsifier{E07、E09均需在观察结果前定义有用增量、独立评价范围与资源边界。差异区间排除所声称幅度的收益，或支持相反方向，才可反对对应范围的主张；仅见训练上涨而迁移差异不显著仍不够。允许重新交互时，也不能假定两路取得相同轨迹。}

\openquestion{第13章的触发讨论与第15章的试跑审计采用外部给定的任务分布。改用自生成课程后，分布相关的结论需要重新检验，预算账还须纳入出题、试做、过滤与材料选择费用。本书尚无这些课程训练的成本实测。}

\section*{习题}
\begin{enumerate}[label=11-\arabic*.,ref=11-\arabic*]
\sloppy
  \item\label{ex:ch11-1} 章首失败案例里，去掉合成器训练与去掉多样性惩罚，AMC23 都从 \cited{62.50} 掉到 \cited{55.00}。
        两行读数相同，原文给出的机制解释却不同。请说明各自的解释，
        并各加一个度量来检验。
  \item\label{ex:ch11-2} \S\ref{sec:ch11-filter} 里，冻结的强合成器带来 \cited{+2.66}，共进化合成器带来 \cited{+5.34}。
        请设计一个能让“出题者绝对强度”与“课程可适应性”分离的对照，
        并说明为什么只在同一个基座上换合成器是不够的。
  \item\label{ex:ch11-3} 分别为任务可行、答案正确、难度合适提出一个可检查指标，指出哪些可由内部代理提供、哪些需要外部依据。
        再用本章二元答案的独立采样设定，计算 $p=0.5$ 时十次回答的多数一致率期望，并与正确率期望比较。若改用模型判断，设计什么对照能发现共同错误？
  \item\label{ex:ch11-4} \S\ref{sec:ch11-reward} 里的格式分可以独立于正确性变化。请把“奖励投机”改写成一句可判定的话
        （出现什么观测就算投机成立），再说明它需要哪个外部量才能被观测到。
  \item\label{ex:ch11-5} 你的自生成课程在训练集上连续十轮上涨。列出三种能产生这条曲线、
        却与“能力提高”无关的解释，并给出各自最便宜的一条检测。
\end{enumerate}

\section*{常见误解}
\begin{itemize}
\sloppy
  \item \textbf{“分数一直在涨就说明课程是对的”。}
        变化课程上的均值可能因题目变容易而上涨；
        固定范围复评或版本与题集的交叉比较有助于辨认方向。不能将这条限制误用于否定论文另行报告的基准评价。
  \item \textbf{“合成器越强越好”。}
        原论文换上更强的冻结合成器只拿到 \cited{+2.66}，
        共进化条件拿到 \cited{+5.34}，但模型身份与训练方式共变，需补充析因对照才能分离。
  \item \textbf{“奖励越密越好，多给中间奖励总没坏处”。}
        奖励本来就应能被策略行为影响；风险在于代理与目标不一致。
        格式项可在正确性不变时提高，应分别观测，不把满足格式要求本身叫作投机。
  \item \textbf{“失败轨迹没有用，扔掉就好”。}
        原文特定蒸馏条件下，保留教师失败轨迹也观察到好处；
        是否分成败取决于设置，不是一条通用规则。
\end{itemize}

\section*{进一步阅读}
\begin{itemize}
\sloppy
  \item \cite{ttcs-2601}：Table 1 的六项均值、Table 2 的合成器冻结对照、
        Table 3 的消融（含本章引用的两行）。读 Table 3 时区分逐部件删除后的效果差与作者提出的机制解释。
  \item \cite{self-challenging-2506}：§1 与 §4.2 的四件套形式，§5.2 Table 1 的
        零样本 \cited{12.0} / 自改进 \cited{23.5} / 蒸馏 \cited{32.2} 三档。
  \item \cite{absolute-zero-2505}：§1、§3 的零数据自博弈与代码执行反馈的构造；
        §6 的 "uh-oh moment" 提醒读者将环境验证与安全治理分别核查。
  \item \cite{agent-lightning-2508}：§2—§3 区分状态与观察，以部分可观测决策过程组织执行、
        把轨迹切成 transition 的具体做法，§3.4 的训练与执行解耦是工程那一节的原形。
  \item 附录~\ref{app:rl}：策略梯度、基线与信用分配的推导，本章只用了结论。
\end{itemize}
